% GENERATED FILE. Edit canonical lessons, then run npm run generate:books.
% canonical-source-hash: fnv1a64:15692ddcbd82fc4b
% canonical-lessons: PA-C314-jammna, PA-C314-niana, PA-C314-palna, PA-C314-pota, PA-C314-dohta

\chapter{A Baby in the House}
\label{ch:pa-a2-a-baby-in-the-house}
\hlchaptermodality{314}

\begin{chapteropening}
\textbf{By the end of this chapter:} \emph{I can give the news of a new baby.}

Complete the last lesson of chapter 314: I can give the news of a new baby.
\end{chapteropening}

\section[\texorpdfstring{\pa{ਜੰਮਣਾ} (jammṇā)}{ਜੰਮਣਾ (jammṇā)}]{\pa{ਜੰਮਣਾ} (jammṇā) — to be born}

\label{lesson:PA-C314-jammna}



\begin{quote}
Before the new one: say the Punjabi for to hand out, then the Punjabi for congratulations.
\end{quote}

\subsection*{You'll want to know: \pa{ਜੰਮਣਾ}}
\textbf{\pa{ਜੰਮਣਾ}} — \emph{jammṇā} — \textquotedblleft{}to be born\textquotedblright{}.

\begin{quote}
\pa{ਮੇਰੀ} \pa{ਭੈਣ} \pa{ਦੀ} \pa{ਧੀ} \pa{ਪਿਛਲੇ} \pa{ਹਫ਼ਤੇ} \pa{ਜੰਮੀ।}
\end{quote}

\begin{quote}
\emph{merī bhaiṇ dī dhī pichhle hafte jammī.}
\end{quote}

\begin{quote}
My sister's daughter was born last week.
\end{quote}

Past (I …): \pa{ਮੈਂ} \pa{ਜਨਵਰੀ} \pa{ਵਿੱਚ} \pa{ਜੰਮਿਆ}, \emph{maiṁ janvarī vich jammiā}.

\begin{cousinweb}
\textbf{\pa{ਜੰਮਣਾ}} (\emph{jammṇā}) is inherited from Sanskrit \emph{janman}, \textquotedblleft{}birth, a creature\textquotedblright{}; CDIAL lists Punjabi \emph{jammṇā}, \textquotedblleft{}to be born\textquotedblright{}, under it (CDIAL 5113).
\end{cousinweb}

\begin{grammarlens}[title={What you've built}]
One more word for saying what happened.
\end{grammarlens}

\subsection*{Your turn}
\begin{itemize}
\raggedright
  \item \emph{Say it:} \emph{jammṇā}
  \item \emph{Say it:} \emph{jammṇā}, once more
  \item \emph{Say it:} \emph{merī bhaiṇ dī dhī pichhle hafte jammī.} — My sister's daughter was born last week.
  \item \emph{Recall:} say the Punjabi for to hand out, then the Punjabi for congratulations, then say \emph{jammṇā} again
\end{itemize}

\begin{tcolorbox}[breakable,colback=teal!4,colframe=teal!35!black,title={Before you move on}]
What does \pa{ਜੰਮਣਾ} mean? (\textquotedblleft{}To be born\textquotedblright{}.) Say it once more.
\end{tcolorbox}
\section[\texorpdfstring{\pa{ਨਿਆਣਾ} (niāṇā)}{ਨਿਆਣਾ (niāṇā)}]{\pa{ਨਿਆਣਾ} (niāṇā) — a small child; a little one}

\label{lesson:PA-C314-niana}



\begin{quote}
Before the new one: say the Punjabi for congratulations, then the Punjabi for to be born.
\end{quote}

\subsection*{You'll want to know: \pa{ਨਿਆਣਾ}}
\textbf{\pa{ਨਿਆਣਾ}} — \emph{niāṇā} — \textquotedblleft{}a small child; a little one\textquotedblright{}.

\begin{quote}
\pa{ਨਿਆਣਾ} \pa{ਰਾਤ} \pa{ਨੂੰ} \pa{ਬਹੁਤ} \pa{ਰੋਇਆ।}
\end{quote}

\begin{quote}
\emph{niāṇā rāt nū̃ bahut roiā.}
\end{quote}

\begin{quote}
The little one cried a lot at night.
\end{quote}

A masculine noun.

\begin{cousinweb}
\textbf{\pa{ਨਿਆਣਾ}} (\emph{niāṇā}) is inherited from Sanskrit \emph{na jānant}, \textquotedblleft{}not knowing\textquotedblright{}; CDIAL lists Punjabi \emph{niāṇā}, \textquotedblleft{}of immature judgement, young, little\textquotedblright{}, under it (CDIAL 6931).
\end{cousinweb}

\begin{grammarlens}[title={What you've built}]
One more word for saying what happened.
\end{grammarlens}

\subsection*{Your turn}
\begin{itemize}
\raggedright
  \item \emph{Say it:} \emph{niāṇā}
  \item \emph{Say it:} \emph{niāṇā}, once more
  \item \emph{Say it:} \emph{niāṇā rāt nū̃ bahut roiā.} — The little one cried a lot at night.
  \item \emph{Recall:} say the Punjabi for congratulations, then the Punjabi for to be born, then say \emph{niāṇā} again
  \item \emph{Write it:} copy \textbf{\pa{ਨਿਆਣਾ}} once, checking the whole word
\end{itemize}

\begin{tcolorbox}[breakable,colback=teal!4,colframe=teal!35!black,title={Before you move on}]
What does \pa{ਨਿਆਣਾ} mean? (\textquotedblleft{}A small child; a little one\textquotedblright{}.) Say it once more.
\end{tcolorbox}
\section[\texorpdfstring{\pa{ਪਾਲਣਾ} (pālṇā)}{ਪਾਲਣਾ (pālṇā)}]{\pa{ਪਾਲਣਾ} (pālṇā) — to bring up}

\label{lesson:PA-C314-palna}



\begin{quote}
Before the new one: say the Punjabi for to be born, then the Punjabi for a small child.
\end{quote}

\subsection*{You'll want to know: \pa{ਪਾਲਣਾ}}
\textbf{\pa{ਪਾਲਣਾ}} — \emph{pālṇā} — \textquotedblleft{}to bring up\textquotedblright{}.

\begin{quote}
\pa{ਅਸੀਂ} \pa{ਬੱਚਿਆਂ} \pa{ਨੂੰ} \pa{ਪਿਆਰ} \pa{ਨਾਲ} \pa{ਪਾਲਿਆ।}
\end{quote}

\begin{quote}
\emph{asīṁ bachchiāṁ nū̃ piār nāl pāliā.}
\end{quote}

\begin{quote}
We brought up the children with love.
\end{quote}

Past (I …): \pa{ਮੈਂ} \pa{ਤਿੰਨ} \pa{ਬੱਚੇ} \pa{ਪਾਲੇ}, \emph{maiṁ tinn bachche pāle}.

\begin{cousinweb}
\textbf{\pa{ਪਾਲਣਾ}} (\emph{pālṇā}) is inherited from Sanskrit \emph{pālayati}, \textquotedblleft{}guards, protects\textquotedblright{}; CDIAL lists Punjabi \emph{pālaṇā}, \textquotedblleft{}to nourish, feed up\textquotedblright{}, under it (CDIAL 8129).
\end{cousinweb}

\begin{grammarlens}[title={What you've built}]
One more word for saying what happened.
\end{grammarlens}

\subsection*{Your turn}
\begin{itemize}
\raggedright
  \item \emph{Say it:} \emph{pālṇā}
  \item \emph{Say it:} \emph{pālṇā}, once more
  \item \emph{Say it:} \emph{asīṁ bachchiāṁ nū̃ piār nāl pāliā.} — We brought up the children with love.
  \item \emph{Recall:} say the Punjabi for to be born, then the Punjabi for a small child, then say \emph{pālṇā} again
  \item \emph{Write it:} copy \textbf{\pa{ਪਾਲਣਾ}} once, checking the whole word
\end{itemize}

\begin{tcolorbox}[breakable,colback=teal!4,colframe=teal!35!black,title={Before you move on}]
What does \pa{ਪਾਲਣਾ} mean? (\textquotedblleft{}To bring up\textquotedblright{}.) Say it once more.
\end{tcolorbox}
\section[\texorpdfstring{\pa{ਪੋਤਾ} (potā)}{ਪੋਤਾ (potā)}]{\pa{ਪੋਤਾ} (potā) — a son's son}

\label{lesson:PA-C314-pota}



\begin{quote}
Before the new one: say the Punjabi for a small child, then the Punjabi for to bring up.
\end{quote}

\subsection*{You'll want to know: \pa{ਪੋਤਾ}}
\textbf{\pa{ਪੋਤਾ}} — \emph{potā} — \textquotedblleft{}a son's son\textquotedblright{}.

\begin{quote}
\pa{ਪਿਛਲੇ} \pa{ਮਹੀਨੇ} \pa{ਦਾਦਾ} \pa{ਜੀ} \pa{ਦਾ} \pa{ਪਹਿਲਾ} \pa{ਪੋਤਾ} \pa{ਜੰਮਿਆ।}
\end{quote}

\begin{quote}
\emph{pichhle mahīne dādā jī dā pahilā potā jammiā.}
\end{quote}

\begin{quote}
Last month my grandfather's first grandson was born.
\end{quote}

A masculine noun.

\begin{cousinweb}
\textbf{\pa{ਪੋਤਾ}} (\emph{potā}) is inherited from Sanskrit \emph{pautra}, \textquotedblleft{}a son's son\textquotedblright{}; CDIAL lists Punjabi \emph{pottā}, \textquotedblleft{}son's son\textquotedblright{}, under it (CDIAL 8416).
\end{cousinweb}

\begin{grammarlens}[title={What you've built}]
One more word for saying what happened.
\end{grammarlens}

\subsection*{Your turn}
\begin{itemize}
\raggedright
  \item \emph{Say it:} \emph{potā}
  \item \emph{Say it:} \emph{potā}, once more
  \item \emph{Say it:} \emph{pichhle mahīne dādā jī dā pahilā potā jammiā.} — Last month my grandfather's first grandson was born.
  \item \emph{Recall:} say the Punjabi for a small child, then the Punjabi for to bring up, then say \emph{potā} again
  \item \emph{Write it:} copy \textbf{\pa{ਪੋਤਾ}} once, checking the whole word
\end{itemize}

\begin{tcolorbox}[breakable,colback=teal!4,colframe=teal!35!black,title={Before you move on}]
What does \pa{ਪੋਤਾ} mean? (\textquotedblleft{}A son's son\textquotedblright{}.) Say it once more.
\end{tcolorbox}
\section[\texorpdfstring{\pa{ਦੋਹਤਾ} (dohtā)}{ਦੋਹਤਾ (dohtā)}]{\pa{ਦੋਹਤਾ} (dohtā) — a daughter's son}

\label{lesson:PA-C314-dohta}



\begin{quote}
Before the new one: say the Punjabi for to bring up, then the Punjabi for a son's son.
\end{quote}

\subsection*{You'll want to know: \pa{ਦੋਹਤਾ}}
\textbf{\pa{ਦੋਹਤਾ}} — \emph{dohtā} — \textquotedblleft{}a daughter's son\textquotedblright{}.

\begin{quote}
\pa{ਨਾਨੀ} \pa{ਆਪਣੇ} \pa{ਦੋਹਤੇ} \pa{ਨਾਲ} \pa{ਬਹੁਤ} \pa{ਖੇਡੀ।}
\end{quote}

\begin{quote}
\emph{nānī āpṇe dohte nāl bahut kheḍī.}
\end{quote}

\begin{quote}
Grandmother played a lot with her daughter's son.
\end{quote}

A masculine noun.

\begin{cousinweb}
\textbf{\pa{ਦੋਹਤਾ}} (\emph{dohtā}) is inherited from Sanskrit \emph{dauhitra}, \textquotedblleft{}a daughter's son\textquotedblright{}; CDIAL lists Punjabi \emph{dohtā}, \textquotedblleft{}daughter's son\textquotedblright{}, under it (CDIAL 6605).
\end{cousinweb}

\begin{grammarlens}[title={What you've built}]
One more word for saying what happened.
\end{grammarlens}

\subsection*{Your turn}
\begin{itemize}
\raggedright
  \item \emph{Say it:} \emph{dohtā}
  \item \emph{Say it:} \emph{dohtā}, once more
  \item \emph{Say it:} \emph{nānī āpṇe dohte nāl bahut kheḍī.} — Grandmother played a lot with her daughter's son.
  \item \emph{Recall:} say the Punjabi for to bring up, then the Punjabi for a son's son, then say \emph{dohtā} again
  \item \emph{Write it:} copy \textbf{\pa{ਦੋਹਤਾ}} once, checking the whole word
\end{itemize}

\begin{tcolorbox}[breakable,colback=teal!4,colframe=teal!35!black,title={Before you move on}]
What does \pa{ਦੋਹਤਾ} mean? (\textquotedblleft{}A daughter's son\textquotedblright{}.) Say it once more.
\end{tcolorbox}
